\documentclass[11pt]{article}
\usepackage{agda}
\usepackage{fontspec}
\usepackage{newunicodechar}
\usepackage{etoolbox}
\usepackage{tikz}
\usetikzlibrary{arrows.meta,positioning}
\setsansfont{Arial Unicode MS}
\usepackage[margin=1in]{geometry}
\usepackage[hidelinks]{hyperref}
\newunicodechar{φ}{\ensuremath{\phi}}

% Match the main Void and Form code layout: compact Agda blocks,
% no extra math indentation, and narrower indentation steps.
\renewcommand{\AgdaCodeStyle}{\fontsize{8}{9}\selectfont}
\setlength{\mathindent}{0pt}
\renewcommand{\AgdaIndentSpace}{\hspace*{0.30em}}
\AtBeginEnvironment{code}{\hfuzz=2pt}

\title{Void Companion: Binary Distinction and the Four-Vertex Closure}
\author{Johannes Michael Wielsch\\
\small Companion note to \emph{Void and Form}\\
\small Project DOI: \href{https://doi.org/10.5281/zenodo.19491035}{10.5281/zenodo.19491035}}
\date{May 2026}

\begin{document}
\maketitle

\begin{abstract}
This companion studies a narrow licensing problem.  A formal system may
say that there are two positions.  What data internal to the system
license that count as two, rather than one position named twice or a
larger carrier with unmentioned points?

The entry pressure is the Leibnizian principle usually called the
identity of indiscernibles: if no expressible difference separates two
alleged positions, the system has no internal ground for holding them
apart.  This principle is not itself proved by Agda.  It motivates the
formal threshold at which the question becomes checkable.

In intensional Martin-L\"of type theory, formalized in Agda under
\texttt{--safe --without-K}, the threshold is represented by a
\AgdaRecord{Distinction} record.  The record contains a carrier
\AgdaField{S}, two distinguished positions \AgdaField{left} and
\AgdaField{right}, a separation witness
\AgdaField{separated} : \AgdaSymbol{\ensuremath{\neg}} (\AgdaField{left}
\AgdaSymbol{\ensuremath{\equiv}} \AgdaField{right}), and a cover
\AgdaField{cover} showing that every carrier element is one of those two
positions.  Thus \AgdaField{separated} is anti-collapse data, and
\AgdaField{cover} is anti-surplus data.

From that record the checked core proves a conditional chain.  Every
\AgdaRecord{Distinction} is boundary-preservingly isomorphic to the
fresh two-point normal form \AgdaDatatype{Two}.  The endomorphism space
\AgdaDatatype{Two} \AgdaSymbol{\ensuremath{\to}} \AgdaDatatype{Two}
has exactly four cases: the two constants, identity, and swap.  Under
the stated representation contract---separation of cases, complete
realisation, and no surplus---the corresponding simple graph closure is
$K_4$.  The closure is then read back to the originating
\AgdaRecord{Distinction} datum.

The result is not a theorem that all formal systems begin this way.  It
is a mechanically checked conditional result inside the stated Agda
fragment.  The broader theory of formulability is the architectural
reading of this chain, not an additional Agda theorem.
\end{abstract}

\section*{Reading key}

The note begins with a small question and keeps returning to it:
\emph{a formal system says that there are two; what inside the system
licenses that claim?}  The answer is not sought in $K_4$, in a
truth-value type, or in a later interpretation.  The answer begins with
the data needed to prevent collapse and hidden surplus.

The pre-formal part introduces the collapse pressure.  The identity of
indiscernibles says, in the form needed here, that two alleged positions
cannot be maintained as two if the system can express no difference
between them.  This is a motivation for the formal datum, not a checked
term of the development.

The formal threshold is crossed only when explicit data are supplied:
carrier, left boundary, right boundary, separation, and cover.  From
that point onward the claims are local and conditional.  Agda checks
what follows from the record; it does not manufacture the pre-formal
reason for asking for the record.

The checked part then follows one route: \AgdaRecord{Distinction} to
\AgdaDatatype{Two}, \AgdaDatatype{Two} to its four endomorphism cases,
the four cases to $K_4$ under the representation contract, and the
closure back to the originating \AgdaRecord{Distinction} datum.  Later
interpretation is kept separate from this kernel.

\begin{center}
\begin{tikzpicture}[
  stage/.style={draw, rounded corners=2pt, align=center, inner xsep=6pt,
    inner ysep=5pt, minimum width=3.1cm},
  flowarrow/.style={-{Latex[length=2mm]}, thick}
]
\node[stage] (question) {licensing problem\\\footnotesize what warrants two?};
\node[stage, right=0.9cm of question] (data) {formal threshold\\\footnotesize separation + cover};
\node[stage, right=0.9cm of data] (checked) {checked chain\\\footnotesize normal form + closure};
\draw[flowarrow] (question) -- (data);
\draw[flowarrow] (data) -- (checked);
\end{tikzpicture}
\end{center}

\section*{Argument map}

\begin{description}
\item[Licensing problem.] The paper starts from the claim of two, not
from a completed two-element object.  The question is what internal data
make such a claim stable.

\item[Collapse pressure.] The identity of indiscernibles supplies the
pressure: without an expressible distinction, there is no internal
warrant for maintaining two alleged positions as two.

\item[Formal threshold.] The \AgdaRecord{Distinction} record supplies
the data used in this note: a carrier, two distinguished positions, a
non-coincidence witness, and an exhaustive cover.

\item[Checked core.] Given the record, the code checks normal form,
four-case endomorphism classification, $K_4$ closure under the stated
representation conditions, and return to the originating datum.

\item[Later layers.] \emph{Form} and the larger \texttt{Void.lagda.tex}
development are later continuations.  They are not premises of the
kernel and do not change the status of the conditional theorem proved
here.
\end{description}

\part*{Part I.\quad The pre-formal threshold}

\section*{The licensing problem}

The pre-formal part has one role: to explain why the formal development
does not begin with an unexamined two-element type.  A relation has two
positions.  A morphism has a source and a target.  A judgement relates a
term and a type.  An equation has two sides, even when it identifies
them.  Each form already uses separability.

The licensing question asks what a system must have in order to keep two
positions apart before further structure is read into them.  Naming two
points is not enough.  If the names can collapse, the datum is not two.
If further points can inhabit the carrier unaccounted, the datum is not
exhaustively two.

The following display is not Agda code and not a theorem of MLTT.  It
is a schematic boundary marker for the passage into the checked part.
The arrows are meta-level dependencies, not implication terms.

\[
\begin{array}{ll}
\mathrm{Level\;0:} &
  \mathcal{V} \notin \mathit{Set}_\omega,
  \qquad \neg\exists x\,(x : \mathcal{V}).
\\[0.6em]
\mathrm{Level\;1:} &
  \mathrm{Sing}(S) :=
  \exists p:S\,\forall x:S\,(x = p).
\\[0.6em]
\mathrm{Collapse:} &
  \mathrm{Sing}(S)
  \Longrightarrow
  \neg\exists a,b:S\,(a \neq b).
\\[0.6em]
\mathrm{Level\;2:} &
  \mathrm{Dist}(S) :=
  \exists a,b:S\,
  \bigl(a \neq b\ \wedge\ \forall x:S\,(x=a \vee x=b)\bigr).
\\[0.6em]
\mathrm{Formal\;threshold:} &
  \mathrm{Dist}(S)
  \Longrightarrow
  \begin{array}[t]{l}
  S \simeq \mathbf{2},\\
  \mathrm{End}(\mathbf{2}) =
    \{c_L,c_R,\mathrm{id},\sigma\},\\
  \mathrm{representation\ closure} = K_4,\\
  K_4 \Longrightarrow \mathrm{Dist}(S).
  \end{array}
\end{array}
\]

\begin{center}
\begin{tikzpicture}[
  node distance=1.0cm,
  state/.style={draw, rounded corners=2pt, align=center, inner xsep=7pt,
    inner ysep=5pt, minimum width=2.8cm},
  pass/.style={-{Latex[length=2mm]}, thick}
]
\node[state] (void) {no carrier\\\footnotesize no typed datum};
\node[state, right=of void] (single) {singleton\\\footnotesize no separation};
\node[state, right=of single] (dist) {Distinction\\\footnotesize no collapse, no surplus};
\node[state, right=of dist] (agda) {Agda fragment\\\footnotesize checked consequences};
\draw[pass] (void) -- (single);
\draw[pass] (single) -- (dist);
\draw[pass] (dist) -- (agda);
\end{tikzpicture}
\end{center}

Level 0 is not a hidden type.  It marks the absence of a domain in which
a term could already be typed.  Level 1 shows why mere inhabitance is
still collapse: if all elements are identified with one point, no
internal separation is available.  Level 2 is the first formal datum used
by the note: two distinguished positions, a separation witness, and a
cover saying that every carrier element is one of those positions.

Only the dependency from \(\mathrm{Dist}(S)\) onward enters Agda.  From
that point, the schematic notation is replaced by checked terms.

\section*{Why two cannot occur without separation}

Before any cardinality is chosen, a sharper question arises.  Why can a
system say that two positions are present when no structure available
inside the system distinguishes them?

The relevant pressure is the identity of indiscernibles, used here in a
modest form: a system cannot maintain a two-count by merely naming two
positions if it has no expressible difference that keeps them apart.
This is the reason the formal datum must include anti-collapse data.

This is why separation is not decoration.  Without a witness that
\AgdaField{left} and \AgdaField{right} do not coincide, the named
positions can collapse.  This is also why cover is not decoration.  Two
named positions do not by themselves rule out a third, unaccounted point
in the carrier.  A stable two-point datum needs both clauses: separation
against collapse, and exhaustion against hidden surplus.

Once a carrier has been chosen, MLTT/Agda gives a concrete execution of
this principle.  The term displayed later as
\AgdaFunction{indiscernibility-implies-equality} is one realization of
the principle, not its source.

\section*{Why two named points are not enough}

The formal answer cannot be merely ``choose two elements''.  In MLTT, a
carrier may contain \AgdaField{left} and \AgdaField{right} and still fail
to be a two-point structure.  It can fail in two ways.

First, the two displayed points may coincide.  The field
\AgdaField{separated} prevents this collapse by requiring a witness of
\AgdaSymbol{\ensuremath{\neg}} (\AgdaField{left}
\AgdaSymbol{\ensuremath{\equiv}} \AgdaField{right}).  Second, the
carrier may contain additional points.  The field \AgdaField{cover}
prevents this surplus by requiring every element to be equal to
\AgdaField{left} or to \AgdaField{right}.  Thus \AgdaField{separated}
is anti-collapse data, and \AgdaField{cover} is exhaustion data.

\section*{Why binary, not unary or ternary}

The binary datum is introduced before any truth-value reading.  At the
pre-formal threshold, the datum must answer two local demands: it must
not collapse to one, and it must not hide further positions.  This
argument is conceptual, not an Agda theorem; the formal results below
do not depend on it.  It is recorded here so that the antecedent of the
machine-checked conditional is not opaque.

\emph{Unary collapses.}  A single occurrence carries no separation.
The schematic predicate $\mathrm{Sing}(S)$ records exactly this: every
element equals the chosen point, and no pair $(a,b)$ with $a\neq b$
can be exhibited.  Whatever appears within $S$ is the same appearance.
Formalisation cannot start here, because relations, morphisms,
judgements, and equations all require two positions before they can be
written down.

\emph{Ternary carries binary data plus surplus.}  Three positions
$a,b,c$ require three pairwise separations $a\neq b$, $b\neq c$,
$a\neq c$.  Each pairwise separation has binary form.  A ternary record
therefore does not remove the binary problem; it adds labels and a cover
over a larger family of positions.  Within the record format studied
here, the ternary case is analysed through its pairwise separations plus
the additional data needed to account for the third label.

\emph{Binary supplies the minimal datum.}  A binary distinction supplies exactly one
separation $a\neq b$ together with the exhaustion clause that nothing
further is present.  No additional structure (order, incidence,
selection) is assumed; any such structure that arises later is obtained
from the four endomorphism cases of the normal form, not added to the
datum.  This is why the formal development below begins with
\AgdaRecord{Distinction} and not with a richer or poorer record.

These observations are later reflected inside the formal development.
A singleton carrier admits no inhabitant of \AgdaRecord{Distinction};
at $n=2$ the n-ary version of the record is mutually convertible with
\AgdaRecord{Distinction}; and higher arities carry pairwise binary
separation witnesses together with a nested binary cover.

\part*{Part II.\quad The formal threshold}

\section{Setting}

Part~I fixed the question that Agda is asked to answer.  The formal
development now begins after the threshold has been specified.  The
setting is ordinary intensional Martin-L\"of type theory as implemented
in Agda.  The code snippets below use the standard library for equality,
sums, products, the empty type, and negation.  This companion records the
minimal spine of the argument; the larger \texttt{Void.lagda.tex}
development rebuilds more infrastructure internally.

The fresh type \AgdaDatatype{Two} has the same shape as
\AgdaDatatype{Bool}, but it is introduced without the truth-value
reading.  In this note, \AgdaDatatype{Two} is used as the normal form of
the \AgdaRecord{Distinction} record, not as a pre-given truth-value
object.

\begin{code}
-- § Safe fragment: no postulates and no Axiom K.
{-# OPTIONS --safe --without-K #-}

-- § Self-contained kernel companion module.
module VoidCompanion where

-- § Standard infrastructure: equality, absurdity, sums, pairs, and negation.
open import Relation.Binary.PropositionalEquality using (_≡_; refl; sym; trans; cong)
open import Data.Empty using (⊥; ⊥-elim)
open import Data.Sum using (_⊎_; inj₁; inj₂)
open import Data.Product using (Σ; _,_)
open import Relation.Nullary using (¬_)
\end{code}

The two-point normal form is the following type.

\begin{code}
-- § The two-point normal form has exactly two constructors.
data Two : Set where
  L R : Two

-- § The two boundary names of Two do not collapse.
L≠R : ¬ (L ≡ R)
L≠R ()

-- § The opposite inequality is obtained by symmetry of equality.
R≠L : ¬ (R ≡ L)
R≠L p = L≠R (sym p)
\end{code}

\section{Distinction}

A binary distinction is not just a type with two named points.  The
record supplies the explicit data that answer the threshold question:
a carrier, two distinguished positions, a proof that those positions do
not coincide, and a cover showing that every carrier element is one of
them.  The \AgdaField{separated} field blocks collapse by
indiscernibility of the two distinguished positions; the
\AgdaField{cover} field blocks hidden
surplus in the carrier.  Together they are the exact hypothesis that
makes the normal-form theorem true.

\begin{code}
-- § Binary distinction: carrier, left/right boundaries, separation, coverage.
record Distinction : Set₁ where
  field
    S         : Set
    left      : S
    right     : S
    separated : ¬ (left ≡ right)
    cover     : (x : S) → (x ≡ left) ⊎ (x ≡ right)
\end{code}

\begin{center}
\begin{tikzpicture}[
  point/.style={circle, draw, thick, minimum size=7mm, inner sep=0pt},
  note/.style={font=\small, align=center},
  link/.style={-{Latex[length=2mm]}, thick}
]
\node[point] (left) at (-2.2,0) {$left$};
\node[point] (right) at (2.2,0) {$right$};
\draw[thick] (left) -- node[above, font=\small] {$separated$} (right);
\node[note] (cover) at (0,-1.15) {$cover(x): x \equiv left$\\or $x \equiv right$};
\draw[link] (cover) -- (left);
\draw[link] (cover) -- (right);
\node[note] at (0,0.85) {carrier $S$ has no hidden surplus};
\end{tikzpicture}
\end{center}

The canonical inhabitant is obtained from \AgdaDatatype{Two} itself.

\begin{code}
-- § The canonical cover of Two is exhaustive by case analysis.
Two-cover : (x : Two) → (x ≡ L) ⊎ (x ≡ R)
Two-cover L = inj₁ refl
Two-cover R = inj₂ refl

-- § The canonical two-point distinction.
Two-distinction : Distinction
Two-distinction = record
  { S         = Two
  ; left      = L
  ; right     = R
  ; separated = L≠R
  ; cover     = Two-cover
  }
\end{code}

\section{Cardinality of distinction}

The earlier prose in Part~I argued that one occurrence collapses,
ternary records retain binary separations plus additional data, and the
binary case supplies the minimal candidate.  The argument was explicitly
marked as conceptual.  At this point, after the
\AgdaRecord{Distinction} record is in place, the same content can be
stated in the precise form available to this framework.  A singleton
carrier admits no inhabitant of \AgdaRecord{Distinction} at all, and
the n-ary generalisation of the record is, at $n=2$, mutually
convertible with the binary case at the displayed carrier and boundary
fields.  Higher arities carry pairwise binary separation witnesses plus
an exhaustive cover that is structurally a nested binary disjunction.
The result is local to this record-based \texttt{--safe --without-K}
fragment.

\begin{code}
-- § Finite arities and finite positions.
open import Data.Nat using (ℕ; zero; suc)
open import Data.Fin using (Fin) renaming (zero to fz; suc to fs)

-- § N-ary distinction template: labelled points, pairwise separation, coverage.
record Distinction-n (n : ℕ) : Set₁ where
  field
    S        : Set
    point    : Fin n → S
    distinct : (i j : Fin n) → ¬ (i ≡ j) → ¬ (point i ≡ point j)
    cover    : (x : S) → Σ (Fin n) (λ i → x ≡ point i)
\end{code}

\emph{Unary collapses formally.}  No singleton carrier hosts a
\AgdaRecord{Distinction}: if every two elements of the carrier are
equal, the boundary points are equal and \AgdaField{separated} is
falsified.

\begin{code}
-- § A carrier whose elements all coincide cannot host a binary distinction.
singleton-no-distinction :
  (d : Distinction) →
  ((x y : Distinction.S d) → x ≡ y) →
  ⊥
singleton-no-distinction d hyp =
  Distinction.separated d (hyp (Distinction.left d) (Distinction.right d))
\end{code}

The same content holds for \AgdaRecord{Distinction-n} at $n=1$: any two
elements of its carrier are identified, so no separation between
distinct positions can be inhabited.

\begin{code}
-- § A unary n-ary distinction carrier collapses every two elements.
distinction-n-1-collapses :
  (d : Distinction-n 1) →
  (x y : Distinction-n.S d) → x ≡ y
distinction-n-1-collapses d x y
  with Distinction-n.cover d x | Distinction-n.cover d y
... | fz , px | fz , py = trans px (sym py)
... | fz , _  | fs () , _
... | fs () , _ | _
\end{code}

The collapse on \AgdaRecord{Distinction-n} \AgdaNumber{1} is one half of
the ``unary kills distinction'' content; on its own, it identifies every
two elements of the carrier to coincide.  The bridge to the original
record is the next step: any \AgdaRecord{Distinction} that injects
into such a carrier is immediately impossible, because the two
separated boundary points become equal under the injection and
\AgdaField{separated} is contradicted.  This makes the joint
statement--a unary carrier hosts no \AgdaRecord{Distinction}--a single
checked term and not a verbal combination of two lemmas.

\begin{code}
-- § No binary distinction injects into a unary n-ary carrier.
no-distinction-on-Dn1 :
  (d1 : Distinction-n 1) →
  (d  : Distinction) →
  (φ  : Distinction.S d → Distinction-n.S d1) →
  ((x y : Distinction.S d) → φ x ≡ φ y → x ≡ y) →
  ⊥
no-distinction-on-Dn1 d1 d φ φ-inj =
  Distinction.separated d
    (φ-inj (Distinction.left d) (Distinction.right d)
      (distinction-n-1-collapses d1
        (φ (Distinction.left d))
        (φ (Distinction.right d))))
\end{code}

\emph{Binary equals the existing record.}  At $n=2$ the n-ary record
is mutually convertible with the original \AgdaRecord{Distinction}.
The two constructions are mechanical and demonstrate that no boundary
information is lost in either direction.  Because the records contain
function fields, the companion does not claim an unqualified record
equality of the conversions; instead it states the recoverable carrier
and boundary fields explicitly.

\begin{code}
-- § Convert a binary distinction into the n-ary template at n=2.
distinction-to-Dn2 : Distinction → Distinction-n 2
distinction-to-Dn2 d = record
  { S        = Distinction.S d
  ; point    = pt
  ; distinct = dist
  ; cover    = cov
  }
  where
    open Distinction d
    pt : Fin 2 → S
    pt fz      = left
    pt (fs fz) = right

    dist : (i j : Fin 2) → ¬ (i ≡ j) → ¬ (pt i ≡ pt j)
    dist fz      fz      h _ = h refl
    dist fz      (fs fz) _ p = separated p
    dist (fs fz) fz      _ p = separated (sym p)
    dist (fs fz) (fs fz) h _ = h refl

    cov : (x : S) → Σ (Fin 2) (λ i → x ≡ pt i)
    cov x with cover x
    ... | inj₁ p = fz , p
    ... | inj₂ p = fs fz , p

  -- § Convert the n=2 template back into the binary distinction record.
Dn2-to-distinction : Distinction-n 2 → Distinction
Dn2-to-distinction d = record
  { S         = Distinction-n.S d
  ; left      = Distinction-n.point d fz
  ; right     = Distinction-n.point d (fs fz)
  ; separated = Distinction-n.distinct d fz (fs fz) λ ()
  ; cover     = cov
  }
  where
    open Distinction-n d
    cov : (x : S) → (x ≡ point fz) ⊎ (x ≡ point (fs fz))
    cov x with cover x
    ... | fz , p    = inj₁ p
    ... | fs fz , p = inj₂ p

-- § Recoverable carrier and boundary fields for the n=2 roundtrips.
record DistinctionDn2Equivalence : Set₁ where
  field
    from-to-carrier   : (d : Distinction) →
      Distinction.S (Dn2-to-distinction (distinction-to-Dn2 d)) ≡ Distinction.S d
    from-to-left      : (d : Distinction) →
      Distinction.left (Dn2-to-distinction (distinction-to-Dn2 d)) ≡ Distinction.left d
    from-to-right     : (d : Distinction) →
      Distinction.right (Dn2-to-distinction (distinction-to-Dn2 d)) ≡ Distinction.right d
    to-from-carrier   : (d : Distinction-n 2) →
      Distinction-n.S (distinction-to-Dn2 (Dn2-to-distinction d)) ≡ Distinction-n.S d
    to-from-point-0   : (d : Distinction-n 2) →
      Distinction-n.point (distinction-to-Dn2 (Dn2-to-distinction d)) fz ≡ Distinction-n.point d fz
    to-from-point-1   : (d : Distinction-n 2) →
      Distinction-n.point (distinction-to-Dn2 (Dn2-to-distinction d)) (fs fz) ≡ Distinction-n.point d (fs fz)

    -- § The n=2 conversions preserve the displayed carrier and boundary fields.
theorem-distinction-Dn2-equivalence : DistinctionDn2Equivalence
theorem-distinction-Dn2-equivalence = record
  { from-to-carrier = λ _ → refl
  ; from-to-left    = λ _ → refl
  ; from-to-right   = λ _ → refl
  ; to-from-carrier = λ _ → refl
  ; to-from-point-0 = λ _ → refl
  ; to-from-point-1 = λ _ → refl
  }
\end{code}

\emph{Higher arity carries only pairwise binary plus a nested
disjunction.}  At any $n$, the n-ary distinction record carries a
binary separation witness for every ordered pair of distinct
positions.  This witness is the \AgdaField{distinct} field itself; no
further structure is needed to extract it.

\begin{code}
-- § Every n-ary distinction supplies binary separation for every distinct pair.
pairwise-binary :
  {n : ℕ} (d : Distinction-n n) →
  (i j : Fin n) → ¬ (i ≡ j) →
  ¬ (Distinction-n.point d i ≡ Distinction-n.point d j)
pairwise-binary = Distinction-n.distinct
\end{code}

At $n=3$, the exhaustive cover of three positions is structurally the
nested binary disjunction $A_0 \uplus (A_1 \uplus A_2)$, where $A_i$
records that an element equals position $i$.  No primitive ternary
disjunction is invoked; the three-case split is two binary splits.

\begin{code}
-- § The three-point cover is a nested binary disjunction.
cover-3-nested :
  (d : Distinction-n 3) (x : Distinction-n.S d) →
  (x ≡ Distinction-n.point d fz)
    ⊎ ((x ≡ Distinction-n.point d (fs fz))
       ⊎ (x ≡ Distinction-n.point d (fs (fs fz))))
cover-3-nested d x with Distinction-n.cover d x
... | fz , p          = inj₁ p
... | fs fz , p       = inj₂ (inj₁ p)
... | fs (fs fz) , p  = inj₂ (inj₂ p)
\end{code}

The two preceding lemmas show only that a ternary record \emph{carries}
binary content.  The stronger claim, that the ternary record is
\emph{nothing more} than this binary content, requires the converse
direction: any package of three pairwise separations together with a
binary-nested cover must reconstruct a \AgdaRecord{Distinction-n}
\AgdaNumber{3}.  The following record names exactly the binary data
extractable from a ternary distinction, and the two conversions below
show the reduction in both directions.

\begin{code}
-- § Binary data package sufficient to reconstruct a ternary distinction.
record Dn3-binary-data : Set₁ where
  field
    S     : Set
    p0    : S
    p1    : S
    p2    : S
    sep01 : ¬ (p0 ≡ p1)
    sep02 : ¬ (p0 ≡ p2)
    sep12 : ¬ (p1 ≡ p2)
    cover : (x : S) → (x ≡ p0) ⊎ ((x ≡ p1) ⊎ (x ≡ p2))

  -- § Extract pairwise binary data from a ternary distinction.
Dn3-to-binary : Distinction-n 3 → Dn3-binary-data
Dn3-to-binary d = record
  { S     = Distinction-n.S d
  ; p0    = Distinction-n.point d fz
  ; p1    = Distinction-n.point d (fs fz)
  ; p2    = Distinction-n.point d (fs (fs fz))
  ; sep01 = Distinction-n.distinct d fz (fs fz) λ ()
  ; sep02 = Distinction-n.distinct d fz (fs (fs fz)) λ ()
  ; sep12 = Distinction-n.distinct d (fs fz) (fs (fs fz)) λ ()
  ; cover = cover-3-nested d
  }

-- § Reconstruct a ternary distinction from pairwise binary data.
binary-to-Dn3 : Dn3-binary-data → Distinction-n 3
binary-to-Dn3 b = record
  { S        = S
  ; point    = pt
  ; distinct = dist
  ; cover    = cov
  }
  where
    open Dn3-binary-data b
    pt : Fin 3 → S
    pt fz             = p0
    pt (fs fz)        = p1
    pt (fs (fs fz))   = p2

    dist : (i j : Fin 3) → ¬ (i ≡ j) → ¬ (pt i ≡ pt j)
    dist fz             fz             h _ = h refl
    dist fz             (fs fz)        _ p = sep01 p
    dist fz             (fs (fs fz))   _ p = sep02 p
    dist (fs fz)        fz             _ p = sep01 (sym p)
    dist (fs fz)        (fs fz)        h _ = h refl
    dist (fs fz)        (fs (fs fz))   _ p = sep12 p
    dist (fs (fs fz))   fz             _ p = sep02 (sym p)
    dist (fs (fs fz))   (fs fz)        _ p = sep12 (sym p)
    dist (fs (fs fz))   (fs (fs fz))   h _ = h refl

    cov : (x : S) → Σ (Fin 3) (λ i → x ≡ pt i)
    cov x with cover x
    ... | inj₁ p          = fz , p
    ... | inj₂ (inj₁ p)   = fs fz , p
    ... | inj₂ (inj₂ p)   = fs (fs fz) , p

-- § Recoverable carrier and boundary fields for the ternary reconstruction.
record Dn3BinaryReconstruction : Set₁ where
  field
    extracted-carrier   : (d : Distinction-n 3) →
      Dn3-binary-data.S (Dn3-to-binary d) ≡ Distinction-n.S d
    extracted-point-0   : (d : Distinction-n 3) →
      Dn3-binary-data.p0 (Dn3-to-binary d) ≡ Distinction-n.point d fz
    extracted-point-1   : (d : Distinction-n 3) →
      Dn3-binary-data.p1 (Dn3-to-binary d) ≡ Distinction-n.point d (fs fz)
    extracted-point-2   : (d : Distinction-n 3) →
      Dn3-binary-data.p2 (Dn3-to-binary d) ≡ Distinction-n.point d (fs (fs fz))
    reconstructed-carrier : (b : Dn3-binary-data) →
      Distinction-n.S (binary-to-Dn3 b) ≡ Dn3-binary-data.S b
    reconstructed-point-0 : (b : Dn3-binary-data) →
      Distinction-n.point (binary-to-Dn3 b) fz ≡ Dn3-binary-data.p0 b
    reconstructed-point-1 : (b : Dn3-binary-data) →
      Distinction-n.point (binary-to-Dn3 b) (fs fz) ≡ Dn3-binary-data.p1 b
    reconstructed-point-2 : (b : Dn3-binary-data) →
      Distinction-n.point (binary-to-Dn3 b) (fs (fs fz)) ≡ Dn3-binary-data.p2 b

    -- § The ternary record is recoverable from pairwise binary data at displayed fields.
theorem-Dn3-binary-reconstruction : Dn3BinaryReconstruction
theorem-Dn3-binary-reconstruction = record
  { extracted-carrier     = λ _ → refl
  ; extracted-point-0     = λ _ → refl
  ; extracted-point-1     = λ _ → refl
  ; extracted-point-2     = λ _ → refl
  ; reconstructed-carrier = λ _ → refl
  ; reconstructed-point-0 = λ _ → refl
  ; reconstructed-point-1 = λ _ → refl
  ; reconstructed-point-2 = λ _ → refl
  }
\end{code}

What these proofs together establish is the formal counterpart of
\emph{Why binary, not unary or ternary} from Part~I, in the precise
sense available to the present framework.  Singleton carriers cannot
host a distinction at all: the lemma
\AgdaFunction{distinction-n-1-collapses} identifies every two elements of
a unary carrier to coincide, and \AgdaFunction{no-distinction-on-Dn1}
closes the bridge by showing that no \AgdaRecord{Distinction} can
inject into such a carrier without contradicting
\AgdaField{separated}.  The original record is exactly the $n=2$
case of an n-ary template (\AgdaFunction{distinction-to-Dn2},
\AgdaFunction{Dn2-to-distinction}), with the carrier and boundary
roundtrips bundled in
\AgdaFunction{theorem-distinction-Dn2-equivalence}.  At $n=3$, the conversions
\AgdaFunction{Dn3-to-binary} and \AgdaFunction{binary-to-Dn3} prove
more than that ternary data carries binary content: they prove that
ternary data is recoverable from binary content, since the binary
package of three pairwise separations and a nested binary cover
reconstructs a \AgdaRecord{Distinction-n} \AgdaNumber{3}.  The
recoverable carrier and boundary fields are collected in
\AgdaFunction{theorem-Dn3-binary-reconstruction}.

The scope of this last claim is precisely the scope of the framework.
The theorem says that, inside this development, every inhabitant of
\AgdaRecord{Distinction-n} \AgdaNumber{3} is recoverable from a binary
package at the displayed carrier and boundary fields; it does not
assert that all function fields are definitionally identical after a
roundtrip, nor does it address alternative frameworks that take a
ternary presentation as primitive.  What is shown is that within the
present record-based, \texttt{--safe --without-K} fragment of MLTT, the
n-ary record at $n=3$ adds no information beyond binary data plus a
labelling.  Binary is therefore the minimal cardinality at which the
\AgdaRecord{Distinction} record is inhabitable in this framework, and
higher arities reduce, in this framework, to binary structure plus a
labelling map.

\section{Normal Form}

An isomorphism of distinctions preserves the two boundary roles and has
inverse maps on carriers.  The normal-form theorem says that every
distinction is isomorphic, in this sense, to the canonical distinction
on \AgdaDatatype{Two}.

\begin{code}
-- § Boundary-preserving isomorphism of distinction presentations.
record DistinctionIso (d e : Distinction) : Set₁ where
  private
    module D = Distinction d
    module E = Distinction e
  field
    to         : D.S → E.S
    from       : E.S → D.S
    to-left    : to D.left ≡ E.left
    to-right   : to D.right ≡ E.right
    from-left  : from E.left ≡ D.left
    from-right : from E.right ≡ D.right
    to-from    : (y : E.S) → to (from y) ≡ y
    from-to    : (x : D.S) → from (to x) ≡ x
\end{code}

The map to \AgdaDatatype{Two} is determined by the cover: an element is sent
to \AgdaInductiveConstructor{L} if it is the left boundary and to
\AgdaInductiveConstructor{R} if it is the right boundary.

\begin{center}
\begin{tikzpicture}[
  box/.style={draw, thick, rounded corners=2pt, minimum width=22mm, minimum height=9mm, align=center, font=\small},
  point/.style={circle, draw, thick, minimum size=7mm, inner sep=0pt, font=\small},
  link/.style={-{Latex[length=2mm]}, thick}
]
\node[box] (s) at (-3,0) {carrier $S$\\with cover};
\node[point] (l) at (0,0.55) {$left$};
\node[point] (r) at (0,-0.55) {$right$};
\node[point] (L) at (3,0.55) {$L$};
\node[point] (R) at (3,-0.55) {$R$};
\draw[link] (s.east) -- (l.west);
\draw[link] (s.east) -- (r.west);
\draw[link] (l) -- node[above, font=\small] {$toTwo$} (L);
\draw[link] (r) -- node[below, font=\small] {$toTwo$} (R);
\draw[link, bend left=22] (L) to node[above, font=\small] {$fromTwo$} (l);
\draw[link, bend right=22] (R) to node[below, font=\small] {$fromTwo$} (r);
\end{tikzpicture}
\end{center}

\begin{code}
-- § The cover determines each carrier point's normal left/right code.
toTwo : (d : Distinction) → Distinction.S d → Two
toTwo d x with Distinction.cover d x
... | inj₁ _ = L
... | inj₂ _ = R

-- § The inverse candidate sends normal points back to the boundary points.
fromTwo : (d : Distinction) → Two → Distinction.S d
fromTwo d L = Distinction.left d
fromTwo d R = Distinction.right d

-- § Left boundary maps to L; the wrong branch contradicts separation.
toTwo-left : (d : Distinction) → toTwo d (Distinction.left d) ≡ L
toTwo-left d with Distinction.cover d (Distinction.left d)
... | inj₁ _ = refl
... | inj₂ p = ⊥-elim (Distinction.separated d p)

-- § Right boundary maps to R; the wrong branch contradicts separation.
toTwo-right : (d : Distinction) → toTwo d (Distinction.right d) ≡ R
toTwo-right d with Distinction.cover d (Distinction.right d)
... | inj₁ p = ⊥-elim (Distinction.separated d (sym p))
... | inj₂ _ = refl

-- § Roundtrip from Two back through the boundary map.
toTwo-fromTwo : (d : Distinction) → (y : Two) → toTwo d (fromTwo d y) ≡ y
toTwo-fromTwo d L = toTwo-left d
toTwo-fromTwo d R = toTwo-right d

-- § Roundtrip from an arbitrary carrier point through Two.
fromTwo-toTwo : (d : Distinction) → (x : Distinction.S d) → fromTwo d (toTwo d x) ≡ x
fromTwo-toTwo d x with Distinction.cover d x
... | inj₁ p = sym p
... | inj₂ p = sym p

-- § Normal form: every distinction is isomorphic to canonical Two.
two-normal-form : (d : Distinction) → DistinctionIso d Two-distinction
two-normal-form d = record
  { to         = toTwo d
  ; from       = fromTwo d
  ; to-left    = toTwo-left d
  ; to-right   = toTwo-right d
  ; from-left  = refl
  ; from-right = refl
  ; to-from    = toTwo-fromTwo d
  ; from-to    = fromTwo-toTwo d
  }

-- § Equality discipline separating definitional and propositional equations.
record NormalFormEqualityDiscipline (d : Distinction) : Set where
  field
    fromTwo-L-definitional : fromTwo d L ≡ Distinction.left d
    fromTwo-R-definitional : fromTwo d R ≡ Distinction.right d
    toTwo-left-propositional : toTwo d (Distinction.left d) ≡ L
    toTwo-right-propositional : toTwo d (Distinction.right d) ≡ R
    toTwo-fromTwo-propositional : (y : Two) → toTwo d (fromTwo d y) ≡ y
    fromTwo-toTwo-propositional : (x : Distinction.S d) → fromTwo d (toTwo d x) ≡ x

-- § Normal-form equality discipline is witnessed by the roundtrip lemmas.
theorem-normal-form-equality-discipline :
  (d : Distinction) → NormalFormEqualityDiscipline d
theorem-normal-form-equality-discipline d = record
  { fromTwo-L-definitional = refl
  ; fromTwo-R-definitional = refl
  ; toTwo-left-propositional = toTwo-left d
  ; toTwo-right-propositional = toTwo-right d
  ; toTwo-fromTwo-propositional = toTwo-fromTwo d
  ; fromTwo-toTwo-propositional = fromTwo-toTwo d
  }
\end{code}

This is the whole normal-form argument.  No choice is made in the map:
the cover field gives only two possible cases, and the separation field
rules out the contradictory boundary cases.  The equality discipline is
explicit.  The map \AgdaFunction{fromTwo} computes by pattern matching,
so its two endpoint equations close by \AgdaInductiveConstructor{refl}.
The equations involving \AgdaFunction{toTwo} are propositional
equalities: they use the supplied cover and, when the wrong boundary is
selected, the separation witness eliminates the impossible branch.

\section{Endomorphisms of Two}

An endomorphism of \AgdaDatatype{Two} is fixed by its two values.  There
are therefore four cases: constant left, constant right, identity, and
swap.

\begin{center}
\begin{tikzpicture}[
  casebox/.style={draw, thick, rounded corners=2pt, minimum width=25mm, minimum height=13mm, align=center, font=\small},
  map/.style={-{Latex[length=1.7mm]}, thick}
]
\node[casebox] (cl) at (-3.2,1.25) {$constL$\\$L\mapsto L,\ R\mapsto L$};
\node[casebox] (cr) at (3.2,1.25) {$constR$\\$L\mapsto R,\ R\mapsto R$};
\node[casebox] (id) at (-3.2,-1.25) {$ident$\\$L\mapsto L,\ R\mapsto R$};
\node[casebox] (sw) at (3.2,-1.25) {$swap$\\$L\mapsto R,\ R\mapsto L$};
\draw[map] (cl) -- (cr);
\draw[map] (id) -- (sw);
\node[font=\small, align=center] at (0,0) {all functions\\$Two \to Two$};
\end{tikzpicture}
\end{center}

\begin{code}
-- § The four possible endomorphism cases of Two.
data EndoCase : Set where
  constL : EndoCase
  constR : EndoCase
  ident  : EndoCase
  swap   : EndoCase

-- § Interpret a case label as an actual endofunction on Two.
interpret : EndoCase → Two → Two
interpret constL _ = L
interpret constR _ = R
interpret ident  x = x
interpret swap   L = R
interpret swap   R = L

-- § Pointwise equality of endofunctions on Two.
Pointwise : (Two → Two) → (Two → Two) → Set
Pointwise f g = (x : Two) → f x ≡ g x
\end{code}

The classifier simply inspects the two values of a function.

\begin{code}
-- § Classify an endofunction by its values at L and R.
classify : (Two → Two) → EndoCase
classify f with f L | f R
... | L | L = constL
... | R | R = constR
... | L | R = ident
... | R | L = swap

-- § Soundness: interpretation of the classified case recovers the function pointwise.
classify-sound : (f : Two → Two) → Pointwise f (interpret (classify f))
classify-sound f L with f L | f R
... | L | L = refl
... | R | R = refl
... | L | R = refl
... | R | L = refl
classify-sound f R with f L | f R
... | L | L = refl
... | R | R = refl
... | L | R = refl
... | R | L = refl
\end{code}

Uniqueness follows by evaluating at \AgdaInductiveConstructor{L} and
\AgdaInductiveConstructor{R}.  Any attempted identification of two
different cases identifies either \AgdaInductiveConstructor{L} with
\AgdaInductiveConstructor{R}, or conversely.

\begin{code}
-- § Distinct case labels cannot have the same interpreted behavior.
case-unique : (c d : EndoCase) → Pointwise (interpret c) (interpret d) → c ≡ d
case-unique constL constL _ = refl
case-unique constL constR p = ⊥-elim (L≠R (p L))
case-unique constL ident  p = ⊥-elim (L≠R (p R))
case-unique constL swap   p = ⊥-elim (L≠R (p L))
case-unique constR constL p = ⊥-elim (R≠L (p L))
case-unique constR constR _ = refl
case-unique constR ident  p = ⊥-elim (R≠L (p L))
case-unique constR swap   p = ⊥-elim (R≠L (p R))
case-unique ident  constL p = ⊥-elim (R≠L (p R))
case-unique ident  constR p = ⊥-elim (L≠R (p L))
case-unique ident  ident  _ = refl
case-unique ident  swap   p = ⊥-elim (L≠R (p L))
case-unique swap   constL p = ⊥-elim (R≠L (p L))
case-unique swap   constR p = ⊥-elim (L≠R (p R))
case-unique swap   ident  p = ⊥-elim (R≠L (p L))
case-unique swap   swap   _ = refl

-- § Classification is unique among all pointwise-correct case labels.
classify-unique :
  (f : Two → Two) → (c : EndoCase) → Pointwise f (interpret c) → classify f ≡ c
classify-unique f c p =
  case-unique (classify f) c
    (λ x → trans (sym (classify-sound f x)) (p x))
\end{code}

\section{Why the Closure is \texorpdfstring{$K_4$}{K4}}

Composition of endomorphisms is total: for any two endomorphism cases,
their composite is again an endomorphism of \AgdaDatatype{Two}, hence
is one of the four cases just classified.

\begin{code}
-- § Composition of endofunctions on Two.
_∘₂_ : (Two → Two) → (Two → Two) → Two → Two
(f ∘₂ g) x = f (g x)

-- § Compose two case labels by classifying the composite function.
composeCase : EndoCase → EndoCase → EndoCase
composeCase c d = classify (interpret c ∘₂ interpret d)

-- § Endomorphisms are closed under composition.
endomorphisms-closed :
  (c d : EndoCase) → Σ EndoCase (λ e → Pointwise (interpret e) (interpret c ∘₂ interpret d))
endomorphisms-closed c d =
  composeCase c d , λ x → sym (classify-sound (interpret c ∘₂ interpret d) x)
\end{code}

The closure is introduced only at this point.  It packages the four
classified endomorphism cases under an explicit representation contract.
Unpacking that contract yields the following three conditions.

\begin{enumerate}
\item \emph{Separation.} Distinct endomorphism cases must remain
  distinct on the carrier.  No identification of cases is allowed.
\item \emph{Realisation.} Every endomorphism case must be represented
  by a vertex of the carrier.
\item \emph{No surplus.} Every vertex of the carrier must arise from
  one of these cases.
\end{enumerate}

These conditions make explicit what ``representation'' means here: each
endomorphism case is carried as a distinct and retrievable element of
the carrier, and every element of the carrier corresponds to such a
case.  Without separation, cases collapse; without realisation, cases
are missing; without the no-surplus condition, additional structure is
introduced.  The contract is stated before it is used.

Under these three conditions, a carrier with fewer than four represented
vertices cannot represent the four cases without identification.  If two
cases are identified, \AgdaFunction{case-unique} shows that the
corresponding functions agree pointwise; but the only way different rows
of the four case table agree is by identifying
\AgdaInductiveConstructor{L} and \AgdaInductiveConstructor{R},
contradicting \AgdaFunction{L≠R}.  Thus separation and realisation
require four represented vertices.

An additional vertex is not licensed by the endomorphism data.  The
composite of any two represented cases is already one of the four
classified cases, by \AgdaFunction{endomorphisms-closed}.  A new vertex
would not be the value of any composition demanded by the case table; it
would violate no surplus.

Finally, total composability is not another assumption.  Any ordered
pair of endomorphisms has a composite, so there is no missing
composability relation between distinct represented cases.  The
underlying simple graph is therefore the complete graph on four
vertices, $K_4$.

\begin{center}
\begin{tikzpicture}[
  vertex/.style={circle, draw, minimum size=1.1cm, inner sep=1pt},
  edge/.style={line width=0.45pt}
]
\node[vertex] (constL) at (90:1.65cm) {$c_L$};
\node[vertex] (ident)  at (0:1.65cm) {$\mathrm{id}$};
\node[vertex] (constR) at (270:1.65cm) {$c_R$};
\node[vertex] (swap)   at (180:1.65cm) {$\sigma$};
\draw[edge] (constL) -- (ident);
\draw[edge] (ident) -- (constR);
\draw[edge] (constR) -- (swap);
\draw[edge] (swap) -- (constL);
\draw[edge] (constL) -- (constR);
\draw[edge] (swap) -- (ident);
\node[align=center, below=2.1cm of constL] {complete simple graph on the four\\[-0.2em]
endomorphism cases of \(\mathbf{2}\)};
\end{tikzpicture}
\end{center}

The following small records package the three conditions.  The record
\AgdaRecord{FaithfulClosure} contains realisation and separation:
\AgdaField{realize} supplies a vertex for each case, and
\AgdaField{reflect-realize} prevents two cases from being identified.
The record \AgdaRecord{MinimalClosure} adds no surplus through
\AgdaField{no-extra}.

These conditions are therefore not independent choices once this notion
of representation has been fixed; they unpack that notion.

\begin{code}
-- § Representation with realisation, reflection, and separation of cases.
record FaithfulClosure : Set₁ where
  field
    V               : Set
    realize         : EndoCase → V
    reflect         : V → EndoCase
    reflect-realize : (c : EndoCase) → reflect (realize c) ≡ c

-- § Minimal closure adds no surplus vertices beyond represented cases.
record MinimalClosure : Set₁ where
  field
    closure  : FaithfulClosure
    no-extra : (v : FaithfulClosure.V closure) →
               Σ EndoCase (λ c → FaithfulClosure.realize closure c ≡ v)

  -- § The canonical vertex carrier is the four-case type itself.
K4Vertex : Set
K4Vertex = EndoCase

  -- § Edges connect distinct cases in the induced simple graph.
K4Edge : K4Vertex → K4Vertex → Set
K4Edge c d = ¬ (c ≡ d)

  -- § Canonical minimal closure of the four endomorphism cases.
canonicalK4 : MinimalClosure
canonicalK4 = record
  { closure = record
      { V               = EndoCase
      ; realize         = λ c → c
      ; reflect         = λ c → c
      ; reflect-realize = λ c → refl
      }
  ; no-extra = λ v → v , refl
  }

-- § No-surplus gives the other roundtrip: realize after reflect returns the vertex.
realize-reflect :
  (m : MinimalClosure) →
  (v : FaithfulClosure.V (MinimalClosure.closure m)) →
  FaithfulClosure.realize (MinimalClosure.closure m)
    (FaithfulClosure.reflect (MinimalClosure.closure m) v)
  ≡ v
realize-reflect m v with MinimalClosure.no-extra m v
... | c , p = trans (sym (cong realize q)) p
  where
    fc = MinimalClosure.closure m
    realize = FaithfulClosure.realize fc
    reflect = FaithfulClosure.reflect fc
    reflect-realize = FaithfulClosure.reflect-realize fc

    q : c ≡ reflect v
    q = trans (sym (reflect-realize c)) (cong reflect p)
\end{code}

The same information can be displayed as a compact $K_4$ presentation.
This is not the large invariant record of the main kernel; it is the
companion-level graph layer induced by any minimal representation of the
four endomorphism cases.  Vertices are the carrier of the closure,
edges connect vertices whose reflected cases are distinct, and the two
roundtrips certify that no case is lost and no vertex is surplus.

\begin{code}
-- § Compact graph-level presentation of the minimal four-case closure.
record K4Presentation : Set₁ where
  field
    vertices              : Set
    edge                  : vertices → vertices → Set
    realize-vertex        : EndoCase → vertices
    reflect-vertex        : vertices → EndoCase
    reflect-realize-vertex : (c : EndoCase) → reflect-vertex (realize-vertex c) ≡ c
    realize-reflect-vertex : (v : vertices) → realize-vertex (reflect-vertex v) ≡ v
    edge-complete         : (v w : vertices) → edge v w ≡ (¬ (reflect-vertex v ≡ reflect-vertex w))
    represented-closure   : MinimalClosure

  -- § Any minimal closure yields a K4 presentation.
minimal-closure-presentation : (m : MinimalClosure) → K4Presentation
minimal-closure-presentation m = record
  { vertices               = FaithfulClosure.V fc
  ; edge                   = λ v w → ¬ (FaithfulClosure.reflect fc v ≡ FaithfulClosure.reflect fc w)
  ; realize-vertex         = FaithfulClosure.realize fc
  ; reflect-vertex         = FaithfulClosure.reflect fc
  ; reflect-realize-vertex = FaithfulClosure.reflect-realize fc
  ; realize-reflect-vertex = realize-reflect m
  ; edge-complete          = λ _ _ → refl
  ; represented-closure    = m
  }
  where
    fc = MinimalClosure.closure m

-- § The canonical K4 presentation is induced by canonicalK4.
canonicalK4Presentation : K4Presentation
canonicalK4Presentation = minimal-closure-presentation canonicalK4
\end{code}

The three record fields are not three independent stipulations.  Each
one is exactly the load-bearing piece of one of the three
representation conditions, and the following lemmas name that load.
The proofs are short because the conditions are exactly the bijection
data; that is the point.

\emph{Separation is load-bearing.}  The field
\AgdaField{reflect-realize} makes \AgdaField{realize} injective.
If two cases were identified on the carrier, applying
\AgdaField{reflect} to the identification would identify the cases
themselves, contradicting the four-case enumeration of
\AgdaDatatype{EndoCase}.

\begin{code}
-- § Separation is load-bearing: realize is injective.
realize-injective :
  (m : MinimalClosure) →
  (c d : EndoCase) →
  FaithfulClosure.realize (MinimalClosure.closure m) c
    ≡ FaithfulClosure.realize (MinimalClosure.closure m) d →
  c ≡ d
realize-injective m c d p =
  trans (sym (FaithfulClosure.reflect-realize (MinimalClosure.closure m) c))
    (trans (cong (FaithfulClosure.reflect (MinimalClosure.closure m)) p)
      (FaithfulClosure.reflect-realize (MinimalClosure.closure m) d))
\end{code}

\emph{Realisation is load-bearing.}  The totality of
\AgdaField{realize} on \AgdaDatatype{EndoCase} is exactly the demand
that every endomorphism case is hosted by some vertex.  Without it,
some case has no representative and the carrier fails to support the
composition table at all.

\begin{code}
-- § Realisation is load-bearing: every case has a vertex.
case-has-vertex :
  (m : MinimalClosure) →
  (c : EndoCase) →
  Σ (FaithfulClosure.V (MinimalClosure.closure m))
    (λ v → FaithfulClosure.realize (MinimalClosure.closure m) c ≡ v)
case-has-vertex m c =
  FaithfulClosure.realize (MinimalClosure.closure m) c , refl
\end{code}

\emph{No surplus is load-bearing.}  The field \AgdaField{no-extra}
forbids vertices that do not arise from any endomorphism case.  Without
it, the carrier may host data that is not licensed by the endomorphism
set, breaking the claim that the closure adds nothing to its datum.

\begin{code}
-- § No surplus is load-bearing: every vertex comes from a case.
vertex-from-case :
  (m : MinimalClosure) →
  (v : FaithfulClosure.V (MinimalClosure.closure m)) →
  Σ EndoCase
    (λ c → FaithfulClosure.realize (MinimalClosure.closure m) c ≡ v)
vertex-from-case m v = MinimalClosure.no-extra m v
\end{code}

Dropping \AgdaField{reflect-realize} re-admits identification of
cases.  Dropping the totality of \AgdaField{realize} re-admits cases
without vertices.  Dropping \AgdaField{no-extra} re-admits vertices
without cases.  In each case, the remaining structure no longer
represents the four-case set in the intended sense.  The three
fields are therefore not chosen because they yield $K_4$; within this
contract they are fixed because dropping any one breaks representation,
and $K_4$ follows.

The companion proves these three load-bearing lemmas against its own
slim records \AgdaRecord{FaithfulClosure} and \AgdaRecord{MinimalClosure}
so that the file remains self-contained and type-checks without the
larger source.  The full development \texttt{Void.lagda.tex} carries
the analogous statements against the richer \AgdaRecord{K4Record} used
there; the present file does not import them, it reproduces them at
the level of detail that this note needs.

This is the sense in which the closure yields $K_4$: the represented
vertex set is the four-case set, and the induced simple graph connects every
distinct pair.  The equations \AgdaField{reflect-realize} and
\AgdaFunction{realize-reflect} show that any carrier satisfying the
three conditions is equivalent to the four-case carrier.

Equivalently, for any \AgdaBound{m} : \AgdaRecord{MinimalClosure}, the
carrier
\[
\AgdaField{FaithfulClosure.V}\,(\AgdaField{MinimalClosure.closure}\,m)
\]
is in bijection with \AgdaDatatype{EndoCase} via \AgdaField{realize}
and \AgdaField{reflect}.  The field \AgdaField{reflect-realize}
proves one composite to be the identity on \AgdaDatatype{EndoCase};
\AgdaFunction{realize-reflect} proves the other composite to be the
identity on the carrier.  If composition on the carrier is read through
this bijection,
\[
v \star_m w =
\AgdaField{realize}_m
  (\AgdaFunction{composeCase}\,(\AgdaField{reflect}_m\,v)
                              (\AgdaField{reflect}_m\,w)),
\]
then reflecting the composite recovers the four-case composition table
of \AgdaDatatype{EndoCase}.  In this precise sense every minimal
closure satisfying separation, realisation, and no surplus is
isomorphic to the canonical four-case carrier.  Thus $K_4$ is not only
minimal; it is unique up to isomorphism under the stated constraints.

\section{Identity of indiscernibles, executed}

The earlier section on why two cannot occur without separation invoked
the pressure of the identity of indiscernibles: if two alleged positions cannot
be distinguished by anything expressible in the system, the system has
no internal ground for holding them apart.  This section does not
introduce that pressure for the first time.  It records only what
becomes possible after a
concrete carrier and equality discipline have already been chosen:
MLTT/Agda executes one precise instance of the principle.  The induction
principle for propositional equality yields it directly.  If every
predicate true of \AgdaBound{a} transfers to \AgdaBound{b}, choose the
predicate \(P\,x = a \equiv x\).  Since \AgdaBound{a} equals itself by
\AgdaInductiveConstructor{refl}, the transfer gives \AgdaBound{a}
\AgdaSymbol{\ensuremath{\equiv}} \AgdaBound{b}.

\begin{code}
-- § Identity of indiscernibles: indistinguishable elements are propositionally equal.
indiscernibility-implies-equality :
  {S : Set} (a b : S) →
  ((P : S → Set) → P a → P b) →
  a ≡ b
indiscernibility-implies-equality a b ind = ind (λ x → a ≡ x) refl
\end{code}

The Agda term is one realisation of the principle, not its source.
Agda does not create the thought that indiscernibles collapse; it
checks this particular execution of it.  Inside the present setting,
the role of \AgdaRecord{Distinction} is now clear: the
\AgdaField{separated} field prevents the two boundary points from
collapsing into one, and the \AgdaField{cover} field prevents a third,
unaccounted position from being hidden in the carrier.

\section{Re-projection to the datum}

The $K_4$ closure does not replace the original distinction; it
presupposes it.  This section makes the dependency explicit at two
levels of strength so that no reader can mistake the closure for an
independent foundation, and so that the trivial direction is not
advertised as more than it is.

\begin{center}
\begin{tikzpicture}[
  item/.style={draw, rounded corners=2pt, align=center, inner xsep=6pt,
    inner ysep=5pt, minimum width=2.7cm},
  route/.style={-{Latex[length=2mm]}, thick}
]
\node[item] (dist) {Distinction\\\footnotesize supplied datum};
\node[item, right=0.7cm of dist] (two) {Two\\\footnotesize normal form};
\node[item, right=0.7cm of two] (endo) {End(Two)\\\footnotesize four cases};
\node[item, right=0.7cm of endo] (kfour) {$K_4$\\\footnotesize representation closure};
\draw[route] (dist) -- (two);
\draw[route] (two) -- (endo);
\draw[route] (endo) -- (kfour);
\draw[route] (kfour.south) to[out=-90,in=-90,looseness=0.85]
  node[below, align=center, font=\footnotesize] {read back\ to the datum} (dist.south);
\end{tikzpicture}
\end{center}

\emph{The structural content as a checked term.}  The bijection
between the carrier of any minimal closure and the four endomorphism
cases is not left to commentary.  It is bundled as a record and
constructed from the lemmas already proved: \AgdaField{reflect-realize}
is one round-trip, \AgdaFunction{realize-reflect} is the other.  This
is the load-bearing structural map.

\begin{code}
-- § Bijection record: two maps and two inverse laws.
record Bijection (A B : Set) : Set where
  field
    to      : A → B
    from    : B → A
    to-from : (b : B) → to (from b) ≡ b
    from-to : (a : A) → from (to a) ≡ a

-- § Every minimal closure carrier is in bijection with EndoCase.
closure-bijection-EndoCase :
  (m : MinimalClosure) →
  Bijection (FaithfulClosure.V (MinimalClosure.closure m)) EndoCase
closure-bijection-EndoCase m = record
  { to      = FaithfulClosure.reflect         (MinimalClosure.closure m)
  ; from    = FaithfulClosure.realize         (MinimalClosure.closure m)
  ; to-from = FaithfulClosure.reflect-realize (MinimalClosure.closure m)
  ; from-to = realize-reflect m
  }
\end{code}

This term is the actual dependency.  Every minimal closure has a
carrier in checked bijection with \AgdaDatatype{EndoCase}; by
\AgdaFunction{two-normal-form}, \AgdaDatatype{EndoCase} is built from
\AgdaDatatype{Two}, the canonical carrier of any binary distinction.
The four-vertex structure is therefore not an independent foundation;
it is, up to the bijection just constructed, the endomorphism set of
the normal form of any binary distinction.

\emph{Canonical re-projection.}  At the level of returning a value of
type \AgdaRecord{Distinction}, every minimal closure re-projects to
the canonical binary distinction.  This is not a carrier-level
isomorphism between four vertices and two points; such an isomorphism
would be false.  The carrier-level bijection is the checked map to
\AgdaDatatype{EndoCase} above.  The re-projection below records the
dependency at the \AgdaRecord{Distinction} level: the four-case closure
is readable only through the binary normal form whose endomorphisms it
represents.

\begin{code}
-- § Reproject a minimal closure to the canonical distinction it presupposes.
k4-presupposes-distinction : MinimalClosure → Distinction
k4-presupposes-distinction _ = Two-distinction

-- § Full checked chain from a distinction to normal form, K4 closure, and return.
record FormalClosureChain (d : Distinction) : Set₁ where
  field
    binary-normal-form        : DistinctionIso d Two-distinction
    equality-discipline       : NormalFormEqualityDiscipline d
    endomorphism-classifier   : (Two → Two) → EndoCase
    classifier-is-sound       : (f : Two → Two) →
      Pointwise f (interpret (endomorphism-classifier f))
    classifier-is-unique      :
      (f : Two → Two) (c : EndoCase) →
      Pointwise f (interpret c) → endomorphism-classifier f ≡ c
    composition-is-closed     :
      (c e : EndoCase) →
      Σ EndoCase (λ r → Pointwise (interpret r) (interpret c ∘₂ interpret e))
    k4-presentation           : K4Presentation
    closure-carrier-bijection : Bijection (K4Presentation.vertices k4-presentation) EndoCase
    reprojected-distinction   : Distinction
    reprojection-is-canonical : reprojected-distinction ≡ Two-distinction

-- § The formal closure chain is witnessed for every binary distinction.
theorem-formal-closure-chain :
  (d : Distinction) → FormalClosureChain d
theorem-formal-closure-chain d = record
  { binary-normal-form        = two-normal-form d
  ; equality-discipline       = theorem-normal-form-equality-discipline d
  ; endomorphism-classifier   = classify
  ; classifier-is-sound       = classify-sound
  ; classifier-is-unique      = classify-unique
  ; composition-is-closed     = endomorphisms-closed
  ; k4-presentation           = canonicalK4Presentation
  ; closure-carrier-bijection = closure-bijection-EndoCase canonicalK4
  ; reprojected-distinction   = k4-presupposes-distinction canonicalK4
  ; reprojection-is-canonical = refl
  }

-- § Same closure chain, starting from the n=2 distinction template.
record Dn2FormalClosureChain (d : Distinction-n 2) : Set₁ where
  field
    n2-equivalence       : DistinctionDn2Equivalence
    left-is-point-0      : Distinction.left (Dn2-to-distinction d) ≡ Distinction-n.point d fz
    right-is-point-1     : Distinction.right (Dn2-to-distinction d) ≡ Distinction-n.point d (fs fz)
    closure-chain        : FormalClosureChain (Dn2-to-distinction d)

-- § The n=2 template closes through the same normal-form and K4 route.
theorem-Dn2-formal-closure-chain :
  (d : Distinction-n 2) → Dn2FormalClosureChain d
theorem-Dn2-formal-closure-chain d = record
  { n2-equivalence     = theorem-distinction-Dn2-equivalence
  ; left-is-point-0    = refl
  ; right-is-point-1   = refl
  ; closure-chain      = theorem-formal-closure-chain (Dn2-to-distinction d)
  }
\end{code}

The full formalisation uses a richer \AgdaRecord{K4Record}.  The
structural dependency carries over without change: the closure carrier
is not an independent foundation, it is, by
\AgdaFunction{closure-bijection-EndoCase}, in checked bijection with
the endomorphism set of the binary normal form.  The record
\AgdaFunction{theorem-formal-closure-chain} bundles the route as one
object, and \AgdaFunction{theorem-Dn2-formal-closure-chain} states the
same route from the $n=2$ template.

\part*{Part III.\quad Frame}

\section{Scope}

The formal result is conditional.  The development never constructs an
inhabitant of \AgdaRecord{Distinction} without supplied data; every
\AgdaRecord{Distinction}-valued construction below the abstract takes a
\AgdaRecord{Distinction} as input or returns the canonical
\AgdaFunction{Two-distinction}.  The code therefore establishes the
consequences of the threshold datum, not the existence of that datum
outside the formal setting.

The conceptual part identifies the dependency that motivates the datum:
relations, morphisms, judgements, and equations use distinguishable
positions before they can be written down.  This is a claim about the
conditions under which formalisation becomes available.  The checked
part begins only after the data are granted.

The machine-checked result is the conditional consequence.  Given a
structure with two distinguishable, exhaustive points, the normal form
is \AgdaDatatype{Two}; given the endomorphisms of that normal form,
representing the four cases with separation, realisation, and no surplus
yields the four-case carrier with complete simple adjacency; and the
closure returns to the datum from which it was built.

The theory of formulability names the architectural reading of this
chain.  It treats the local question answered here--what licenses a
system to keep two points apart--as the first inspectable threshold of
formal description.  The checked kernel explains what follows from the
stated datum; the larger \texttt{Void.lagda.tex} development asks how
far the required infrastructure can be rebuilt internally.

\section{Placement}

This section locates the contribution so that the formal result is read
at the right level.

\emph{Leibniz.}  The pressure behind the entry question is not invented
by the companion.  It is the thought behind the identity of
indiscernibles: if no
expressible difference separates two alleged positions, the system has
no internal ground for holding them apart.  The Agda term
\AgdaFunction{indiscernibility-implies-equality} is one execution of
this pressure after a carrier has been chosen, not the source of the
pressure itself.

\emph{Lawvere.}  ``Adjointness in Foundations'' (1969) treats truth
and falsity as adjoint to the act of distinguishing.  The present note
operates one level below: before truth-values are assigned, before
propositions are evaluated.  \AgdaDatatype{Two} is a fresh carrier,
deliberately unread as $\{\bot,\top\}$, so that the closure result
depends on separation alone and not on a prior reading of the two
points as truth-values.

\emph{Univalent foundations.}  \AgdaDatatype{Two} is an h-set, and the
development uses \texttt{--without-K}.  No higher-dimensional
structure is assumed or invoked.  The note is therefore compatible
with both intensional Martin-Löf type theory and a univalent reading;
nothing in the closure result depends on the choice of identity
discipline.

The contribution is a machine-checked conditional statement at a precise
threshold: given explicit separation and cover, the normal form is
\AgdaDatatype{Two}; its endomorphisms have four cases; under the stated
representation conditions the corresponding closure is $K_4$, unique up
to isomorphism; and that closure points back to the data from which it
was built.

\section{Conclusion and Outlook}

The opening question was: a formal system says that there are two; what
inside the system licenses that claim?  The answer delivered by the
checked part of this note is deliberately conditional.  Once a carrier,
two boundary points, a proof of non-coincidence, and an exhaustive cover
are supplied, the datum has canonical two-point normal form.  Its
endomorphisms classify into exactly four cases.  Under separation,
complete realisation, and no surplus, the representation of those four
cases is the $K_4$ closure.  That closure does not replace the original
datum; it points back to it.

This closes the companion's formal question.  The note shows what
follows when the threshold data are made explicit.  The result is a
small closed argument: the beginning asks what licenses two; the end
answers that two is licensed, in this formal fragment, by separation
against collapse and cover against surplus.

The outlook therefore points to \texttt{Void.lagda.tex}, not to material
outside the formal build.  The larger file continues the same discipline:
keep the formal infrastructure inside the development as far as Agda
permits.  Apart from the primitive universe machinery needed by Agda,
the later construction does not import ready-made equality, arithmetic,
rationals, or real analysis as background.  It rebuilds the route from
natural numbers to integers, from integers to rationals, and from
rational distance and order toward Cauchy-style real structure.

In that sense the next question is not what the closure means outside
the formal development.  The next question is how much of the machinery
usually taken for granted can be reconstructed after the binary
threshold has been fixed.  The companion isolates the first closed
kernel; \texttt{Void.lagda.tex} tests how far that kernel can carry a
self-contained formal build.

\section{Navigation into Void}

The corresponding self-contained development is in the repository
\href{https://github.com/de-johannes/Void-and-Form}{de-johannes/Void-and-Form}.
The main file is \texttt{Void.lagda.tex}.  The route through the larger
formal development is:

\begin{itemize}
\item \emph{Void as boundary}: \emph{What Void Does Not Name},
  \emph{Why Formal Systems Cannot Begin Themselves}, and
  \emph{Representation Requires A Cut}.
\item \emph{The binary threshold}: \emph{Why The First Cut Is Binary},
  \emph{The Formal Threshold}, and the record \texttt{Distinction}.
\item \emph{Normal form and classification}: \emph{Elimination and
  Classification}, \emph{Two Normal Form}, and the endomorphism
  classification of \texttt{Two -> Two}.
\item \emph{Four-vertex closure}: \emph{Graph Presentation},
  \emph{The K4 Invariant Record}, \emph{K4 Representation Theory},
  and \texttt{K4Record-is-canonical}.
\item \emph{Return to the datum}: \texttt{record-presupposes-distinction},
  the formal witness that the surviving closure entails the originating
  distinction.
\item \emph{Internal arithmetic}: the construction of natural numbers,
  integer normal forms, rational numbers, rational equality, order, and
  distance inside the file rather than by standard-library import.
\item \emph{Cauchy route}: the rational setoid and order laws, triangle
  inequality infrastructure, and Cauchy-style convergence material that
  begin the passage from rational numbers toward real structure.
\end{itemize}

\end{document}
